\section{Related literature}\label{sec:lit}

\paragraph{Order imbalance and order-book predictability.}
\citet{chordia2002} show that aggregate order imbalance predicts market returns, and \citet{chordia2004} that stock-level imbalance predicts next-day individual returns, with the effect reversing at longer lags. \citet{cont2014} relate price changes to order-flow imbalance at the best quotes, and \citet{cont2023} extend this to multiple book levels and cross-impact. \citet{kolm2023} find that deep networks fed with stationary order-flow inputs forecast returns at multiple short horizons. \citet{bogousslavsky2023} show that the volatility of order imbalance drives liquidity and is priced in the cross-section. \citet{akyildirim2021} show that order-imbalance analytics sold by Borsa Istanbul predict intraday excess returns. Our features also draw on measures of price impact \citep{kyle1985}, illiquidity \citep{amihud2002}, flow toxicity \citep{easley2012}, range-based volatility \citep{parkinson1980} and realized skewness \citep{amaya2015}; see \citet{bouchaud2018} for a synthesis. We differ in horizon (ten trading days rather than seconds to one day), in cross-sectional evaluation, and in our focus on whether a subscriber facing realistic frictions can monetize the signal.

\paragraph{From anomaly returns to net returns.}
\citet{novymarx2016} show that trading costs eliminate many anomalies, especially high-turnover ones, and that buy/hold spreads mitigate costs. \citet{garleanu2013} show that optimal trading with predictable returns and quadratic costs moves partially toward a target portfolio, which motivates our partial-adjustment books. \citet{muravyev2025} find that the average long--short anomaly return comes from the short leg and falls to roughly zero once borrow fees are charged. \citet{jacobs1998} and \citet{grinold2000b} show that the relative size of long and short exposures is a design choice with first-order consequences, a point we return to in Section~\ref{sec:mapping}.

\paragraph{Overfitting, multiple testing, leakage and long horizons.}
The deflated Sharpe ratio \citep{bailey2014dsr}, the probability of backtest overfitting \citep{bailey2017pbo}, minimum track-record length \citep{bailey2012sharpe} and the statistics of Sharpe ratios \citep{lo2002} address selection among backtests. Reality-check and step-wise procedures \citep{sullivan1999,white2000,hansen2005,romano2005}, false-discovery control \citep{benjamini1995,benjamini2001} and recent work on thresholds and replication \citep{harvey2016,harvey2020,chordia2020,hou2020,giglio2021} address multiplicity; \citet{novymarx2015} shows that combining several signals chosen in-sample inflates $t$-statistics far beyond conventional levels. Purging and embargoing overlapping labels follow \citet{lopezdeprado2018}; \citet{kaufman2012} formalize leakage, and \citet{zhang2026} show that single protocol deviations can manufacture apparent alpha. Overlapping long-horizon outcomes reduce the effective sample and can make conventional standard errors over-confident \citep{hjalmarsson2011,boudoukh2019}. Our protocol borrows from pre-registration \citep{nosek2018} and equivalence testing \citep{lakens2017,lakens2018}, and our reporting follows the checklist logic of \citet{khushi2026}.

\paragraph{Machine learning and foundation models.}
\citet{gu2020} document gains from machine learning in cross-sectional return prediction. Evidence for pretrained time-series foundation models is mixed: \citet{rahimikia2025} and \citet{alonso2026} find small or inconsistent gains for off-the-shelf models such as TimesFM \citep{das2024}. We report a negative transfer result consistent with that evidence.

\paragraph{The Russian equity market.}
Studies of the Russian market have focused on sanctions \citep{ankudinov2017,dovbnya2020} and, recently, news sentiment \citep{makeeva2026}, which finds sentiment to be mainly a reactive indicator. We are not aware of published work on the cross-sectional content of AlgoPack order-flow data or of real-time evaluations that account for its regenerated history. One of our portfolio modules replicates the higher-moment VAR-ICA-MGARCH approach of \citet{mikhaylov2026}; see also \citet{lassance2022}.
